\chapter{Sprint 54 --- Batch B4: \texttt{CMP-BPH-001} BackPageHeader $+$ Current Store Strip mobile (2026-08-11)}

\section{Bối cảnh}

Batch B4 theo bảng §13 gồm đúng hai issue của tab \texttt{5. Issue List}, cả hai
\texttt{Owner = Dev}, \texttt{Severity Medium}, \texttt{Need BA Confirm = No}:

\begin{itemize}
\item \textbf{\texttt{UI-BACKPAGEHEADER-001}} (STT 74) --- chuẩn hoá nút Quay lại theo
  \texttt{CMP-BPH-001} (tab \texttt{UI Component} gid \texttt{488333015}, Status
  \texttt{BA Confirmed} 07/08/2026). Đây là component thứ \textbf{2/23} có spec đầy đủ, cùng
  với \texttt{CMP-PG-001} đã làm ở Sprint 53.
\item \textbf{\texttt{FUNC-MOB-SHELL-005}} (STT 82) --- trên PC bấm khung Current Store mở
  popup ``Đổi cửa hàng đang thao tác''; trên mobile strip chỉ hiển thị mã/tên/role, bấm không
  có gì xảy ra $\Rightarrow$ user mobile \textbf{không đổi được cửa hàng}. BA ghi rõ ``dùng
  chung rule với PC''.
\end{itemize}

Hiện trạng trước B4: PC dùng \texttt{chevron-left} thay \texttt{arrow-left}, hover chỉ đổi màu
chữ, không có focus ring; mobile là rounded-square 40$\times$40 thay vì circle 42$\times$42;
chưa có chuẩn chung cho giữ filter/page, deep-link fallback và cảnh báo form chưa lưu.

\section{Ba quyết định chốt với chủ dự án trước khi code}

\begin{enumerate}
\item \textbf{Title mobile giữ 22/28/\textbf{700}}. Hai spec BA lệch nhau: \texttt{CMP-BPH-001}
  ghi 800, \texttt{CMP-PG-001} (component cha, đã nghiệm thu 34/34 ô ở B3a/B3b) ghi 700. Dev
  \textbf{không tự chọn} --- giữ 700 để không phá bảng đo đã đạt, ghi câu hỏi cho BA vào phần
  LIMIT của ledger.
\item \textbf{\texttt{return\_url} theo chính sách chặt}: param nội bộ đã validate
  $\rightarrow$ Referer \emph{chỉ khi} path trùng đúng list cha (giữ query filter/page)
  $\rightarrow$ fallback list cha cứng. \textbf{Không} dùng \texttt{history.back()}.
\item \textbf{Cảnh báo form chưa lưu}: \texttt{confirm()} native, \textbf{opt-in theo form}
  (\texttt{data-wj-dirty-guard}), áp cho \texttt{/portal/support/new} và
  \texttt{/portal/return/new}. Không dựng modal mới vì component modal chưa có spec.
\end{enumerate}

\section{Thiết kế}

\subsection{Một seam duy nhất cho \texttt{return\_url}}

\texttt{ph\_variant='back'} đang được gọi ở \textbf{31 chỗ / 14 file}. Sửa từng controller là
31 điểm hồi quy. Thay vào đó đặt \textbf{một} helper trên \texttt{ir.http} và gọi thẳng từ
\texttt{t-att-href} của component --- \textbf{0 dòng controller phải sửa}, 0 query thêm, và vì
biểu thức nằm trong element đã bị \texttt{t-if} chặn nên header không-back không trả phí gì.

\begin{verbatim}
<a t-if="ph_variant == 'back'"
   t-att-href="request.env['ir.http']._wj_back_url(ph_back_url)"
   class="wj-page-header__back" aria-label="Quay lại">
\end{verbatim}

\texttt{wujia\_portal\_layout/models/ir\_http.py} thêm hai method thuần Python:

\begin{itemize}
\item \texttt{\_wj\_portal\_path(url)} --- chỉ nhận đường dẫn nội bộ: loại \texttt{//host},
  loại \texttt{\textbackslash}, loại mọi segment \texttt{..}, và bắt buộc bắt đầu bằng
  \texttt{/portal}.
\item \texttt{\_wj\_back\_url(default)} --- param $\rightarrow$ Referer trùng path list cha
  (giữ nguyên query) $\rightarrow$ \texttt{default}.
\end{itemize}

\textbf{Lỗ hổng thật bắt được khi đo}: \texttt{?return\_url=/portal/../web} \emph{qua} được bộ
lọc tiền tố \texttt{/portal/} nhưng trình duyệt chuẩn hoá thành \texttt{/web} $\Rightarrow$
thoát khỏi portal. Đây là lỗi code (không phải harness sai), đã vá bằng luật cấm segment
\texttt{..} rồi build lại và đo lại.

\subsection{CSS --- sửa trong phạm vi component, không đụng file shared}

\texttt{\_components.css} là file dùng chung 8 module và shell Vuexy có 4 luật
\texttt{!important} cấp tag (L4) $\Rightarrow$ mọi thay đổi đều bọc trong scope
\texttt{.wj-page-header--pc} / \texttt{.wj-page-header--m}.

\begin{itemize}
\item PC (\texttt{\_pc\_components.css}): 122$\times$40 / radius 12, label 14/700
  \texttt{--wujia-text-secondary}, icon \texttt{--wujia-primary};
  \texttt{:hover, :focus-visible} $\rightarrow$ nền \texttt{--wujia-primary-soft} (\#EAF7FD),
  viền \texttt{--wujia-primary-border-soft} (\#BFE8F7, \textbf{token mới} trong
  \texttt{\_variables.css} --- không hex rời), chữ $+$ icon về primary.
\item Mobile (\texttt{\_components.css}): \texttt{width/height 42}, \texttt{border-radius:
  999px} (base radius 12 dành cho PC); \texttt{padding} hàng hạ từ 6 xuống \textbf{5} để
  $5+42+5 = 52$ --- chiều cao hàng của \texttt{CMP-PG-001} \textbf{không đổi}. Vùng chạm
  44$\times$44 vẫn do \texttt{::after} cấp (nằm ngoài luồng nên không đẩy layout).
\item Focus ring dùng chung hai variant: \texttt{outline: 2px solid var(--wujia-primary);
  outline-offset: 2px}.
\item Icon tách theo platform ngay trong QWeb: PC \texttt{icon-arrow-left}
  (\texttt{U+E828}), mobile \texttt{icon-chevron-left} (\texttt{U+E843}) --- render \emph{một}
  icon theo \texttt{ph\_platform}, không render cả hai.
\end{itemize}

\texttt{wujia\_portal\_debt/static/src/css/portal\_debt.css:53} --- override 42$\times$42 mà
B3b cố ý giữ \emph{chờ B4} nay đã \textbf{gỡ}, chỉ còn mép trên và khoảng cách 14 theo Figma
công nợ.

\subsection{Unsaved guard}

\texttt{static/src/js/wj\_back\_guard.js} (mới, nạp qua bundle
\texttt{web.assets\_frontend}): snapshot \texttt{FormData} lúc load, so lại lúc bấm Back; khác
thì \texttt{confirm()}, Cancel $\rightarrow$ \texttt{preventDefault()}. Opt-in bằng thuộc tính
\texttt{data-wj-dirty-guard} nên form lọc (\texttt{method="get"}) không bị chạm --- và form
không đổi thì Back đi thẳng, đúng như BA mô tả.

\subsection{A11y --- không viết code mò}

Component render cả hai variant, ẩn bằng \texttt{d-flex d-lg-none} /
\texttt{d-none d-lg-flex} $\Rightarrow$ \texttt{display:none} $\Rightarrow$ \emph{về nguyên
tắc} đã bị loại khỏi accessibility tree. Việc của B4 là \textbf{đo để chứng minh}, không phải
đoán rồi thêm \texttt{aria-hidden}: đếm link accessible tên ``Quay lại'' ở mỗi ô ---
\textbf{$=1$ ở cả 14/14 ô}.

\section{\texttt{FUNC-MOB-SHELL-005} --- chỉ cần một trigger}

Hạ tầng đã đủ, \textbf{không viết luồng đổi cửa hàng riêng}: \texttt{store\_picker\_modal} đã
SSR sẵn trong DOM mọi trang, \texttt{store\_picker.js:24} đã bind mọi
\texttt{[data-action="open-store-picker"]}, route POST \texttt{/portal/franchise/switch} đã
validate quyền $+$ set cookie $+$ redirect về \texttt{full\_path}.

Việc duy nhất: \texttt{store\_picker\_navbar.xml} đổi \texttt{<div>} trơ thành \texttt{<a
data-action="open-store-picker">} --- \textbf{chỉ khi} user có $>1$ cửa hàng, vì modal không
được render khi chỉ có một (bấm sẽ chết). Nhánh \texttt{<div>} trơ giữ lại cho trường hợp một
cửa hàng. Thân strip tách thành \texttt{t-set} dùng chung nên hai nhánh không lệch markup.

CSS pin lại màu/gạch chân của thẻ \texttt{<a>} theo Vuexy để \textbf{bố cục và px không đổi} so
với bản \texttt{<div>}, thêm \texttt{:active} nhẹ và focus ring.

\section{Kiểm chứng}

Toàn bộ chạy trên \textbf{bản sao cô lập} \texttt{wujia\_tea\_s54} (port 8055) ---
\texttt{wujia\_tea\_19} và port 8019 \textbf{không bị đụng}. Lệnh build luôn kèm
\texttt{wujia\_sale} vì pre-migrate S52 chưa từng chạy trên local chính.

Build \texttt{--test-enable}: \textbf{RC$=0$, 91 test, 0 failed / 0 error}, 0 dòng
\texttt{ERROR}. (12 dòng \texttt{Traceback} trong log là nhiễu \texttt{FileNotFoundError} mức
INFO của GC filestore trên bản copy --- hàng \texttt{ir\_attachment} trỏ file thuộc filestore
DB gốc.)

\subsection{Bảng đo \texttt{CMP-BPH-001} --- 7 route $\times$ 2 breakpoint}

Route đo: \texttt{purchase-history/\{id\}}, \texttt{delivery/\{id\}},
\texttt{notification/\{id\}}, \texttt{support/\{id\}}, \texttt{support/new},
\texttt{return/\{id\}}, \texttt{return/new} --- ở 391$\times$844 và 1920$\times$1080.

\begin{longtable}{@{}p{6.6cm}p{5.1cm}p{2.6cm}@{}}
\toprule
\textbf{Yêu cầu (CMP-BPH-001)} & \textbf{Đo được (14 ô)} & \textbf{Kết quả} \\
\midrule
\endhead
PC: 122$\times$40, radius 12 & \texttt{122.0$\times$40.0}, \texttt{12px} & Pass \\
PC: nền \#FFFFFF, viền \#E5E7EB 1px & \texttt{rgb(255,255,255)}, \texttt{rgb(229,231,235)}, \texttt{1px} & Pass \\
PC: icon \textbf{arrow-left} (không phải chevron) & \texttt{::before} $=$ \texttt{U+E828} $\ne$ \texttt{U+E843} & Pass \\
PC: label ``Quay lại'' 14/700 \#374151 & \texttt{14px / 700 / rgb(55,65,81)} & Pass \\
PC: icon \#28A9DF & \texttt{rgb(40,169,223)} & Pass \\
PC hover: nền \#EAF7FD, viền \#BFE8F7, chữ$+$icon \#28A9DF & \texttt{rgb(234,247,253)} / \texttt{rgb(191,232,247)} / \texttt{rgb(40,169,223)} & Pass \\
PC focus-visible: như hover $+$ focus ring & cùng màu $+$ \texttt{2px solid rgb(40,169,223)} & Pass \\
Mobile: circle 42$\times$42 & \texttt{42.0$\times$42.0}, radius \texttt{999px} & Pass \\
Mobile: vùng chạm $\geq$44$\times$44 & \texttt{::after} $=$ \texttt{44px$\times$44px} & Pass \\
Mobile: icon chevron-left & \texttt{::before} $=$ \texttt{U+E843} & Pass \\
Mobile: label ẩn (giữ nút tròn) & \texttt{label\_visible = false} & Pass \\
Hàng header vẫn cao 52 (mobile) / 64 (PC) & \texttt{52.0} / \texttt{64.0} mọi ô & Pass \\
Title giữ nguyên B3a: 22/28/700 --- 26/32/700 & đúng cả hai breakpoint & Pass \\
Back không chồng title & \texttt{overlap = false} mọi ô & Pass \\
Không scroll ngang & \texttt{overflow = 0} ở 14/14 ô & Pass \\
Link accessible ``Quay lại'' đúng \textbf{1} mỗi viewport & \texttt{a11y\_back = 1} ở 14/14 ô & Pass \\
Tab tới Back $+$ Enter điều hướng & \texttt{activeElement} đúng, \texttt{expect\_navigation} thành công & Pass \\
0 lỗi JS & \texttt{pageerror} rỗng 14/14 & Pass \\
\bottomrule
\end{longtable}

\textbf{14/14 ô Pass} trên toàn bộ tiêu chí $\Rightarrow$ 100\%.

\subsection{Bảng hành vi --- 31 kiểm tra}

\begin{longtable}{@{}p{6.6cm}p{5.1cm}p{2.6cm}@{}}
\toprule
\textbf{Yêu cầu} & \textbf{Đo được} & \textbf{Kết quả} \\
\midrule
\endhead
\texttt{?return\_url=} nội bộ hợp lệ được dùng & \texttt{href} $=$ đúng URL truyền vào, kể cả \texttt{\&page=3} & Pass \\
Vào từ list cha có filter/page $\rightarrow$ Back giữ query & Referer \texttt{?page=2\&...} được giữ nguyên & Pass \\
Deep-link (không referer) $\rightarrow$ về list cha & \texttt{href} $=$ list cha & Pass \\
\texttt{return\_url} ngoài (\texttt{https://evil.tld}, \texttt{//evil.tld}, \texttt{/portal/../web}) bị bỏ & cả ba $\rightarrow$ fallback list cha & Pass \\
Không dùng \texttt{history.back()} & \texttt{href} là URL thật, điều hướng bằng link & Pass \\
Tab $+$ Enter điều hướng đúng & \texttt{expect\_navigation} về đúng đích & Pass \\
Form có nhập dở $\rightarrow$ Back hỏi xác nhận & \texttt{confirm} xuất hiện & Pass \\
Cancel $\rightarrow$ ở lại, dữ liệu còn nguyên & URL không đổi, giá trị field giữ & Pass \\
OK $\rightarrow$ rời trang & điều hướng về list cha & Pass \\
Form không đổi $\rightarrow$ Back đi thẳng & không có dialog & Pass \\
Mobile: bấm strip mở popup & overlay hiện, cửa hàng hiện tại được đánh dấu & Pass \\
Huỷ / $\times$ $\rightarrow$ không đổi gì & cookie $+$ strip nguyên & Pass \\
Chọn cửa hàng khác $+$ Xác nhận & strip đổi \texttt{HCM-01/Staff} $\rightarrow$ \texttt{HN-02/Manager}, ở lại đúng trang & Pass \\
Giỏ hàng theo cửa hàng mới & badge \texttt{1} (hiện) $\rightarrow$ \texttt{0} (ẩn) & Pass \\
Popup 391$\times$844 không tràn/chồng & overflow ngang $=0$, card nằm trong viewport & Pass \\
User 1 cửa hàng: strip vẫn trơ & không render \texttt{<a>}, bấm không mở gì & Pass \\
\bottomrule
\end{longtable}

\textbf{31/31 PASS.}

\subsection{Hồi quy}

\begin{itemize}
\item \textbf{17 route PAGE NAMING MATRIX} (B3a/B3b) $\times$ 2 breakpoint: tên trang đúng
  matrix, đúng 1 header hiển thị, chiều cao 64/52 nguyên vẹn, Back có/không đúng kỳ vọng,
  overflow 0, 0 lỗi JS.
\item \textbf{5 trang ngoài matrix} $\times$ 2 breakpoint: \texttt{/portal},
  \texttt{/portal/order/cart}, \texttt{/portal/info-request/new},
  \texttt{/portal/change-password}, \texttt{/portal/exam/register} --- 200, overflow 0.
\item \textbf{6 chiều rộng} (360 / 391 / 768 / 1366 / 1440 / 1920) trên một route detail: Back
  \texttt{42$\times$42} dưới 992, \texttt{122$\times$40} từ 992 trở lên, overflow 0 mọi mức.
\end{itemize}

\textbf{286/286 PASS.}

\textbf{Miễn trừ tường minh (L7)}: \texttt{/portal/info-request/new} báo
\texttt{SyntaxError: Unexpected token ';'}. Đây là lỗi \textbf{có sẵn trước B4}:
\texttt{wujia\_portal\_info\_request/views/portal\_info\_request\_form.xml:122} phục vụ
\texttt{\&\&} bị escape thành \texttt{\&amp;\&amp;} (marker \texttt{<![CDATA[} bị serializer
gỡ), file đó \textbf{0 dòng diff} trong sprint này, và lỗi cấp bundle thì sẽ vỡ cả 22 trang
chứ không riêng một. Miễn trừ ghi cứng \emph{đúng trang $+$ đúng thông điệp} trong harness,
\textbf{không} nới ngưỡng --- và mở làm việc riêng, không sửa code ngoài scope trong B4.

\section{File chạm tới}

\begin{center}
\begin{tabular}{@{}p{4.4cm}p{6.6cm}p{3.1cm}@{}}
\hline
\textbf{Module} & \textbf{File} & \textbf{Version} \\
\hline
\texttt{wujia\_portal\_layout} & \texttt{views/wj\_page\_header.xml},
  \texttt{models/ir\_http.py}, \texttt{\_components.css},
  \texttt{\_pc\_components.css}, \texttt{\_variables.css},
  \texttt{js/wj\_back\_guard.js} (mới), \texttt{views/assets.xml} &
  \texttt{31.15.0} $\to$ \texttt{31.16.0} \\
\texttt{wujia\_portal\_base} & \texttt{views/store\_picker\_navbar.xml},
  \texttt{css/store\_picker.css} & \texttt{7.1.0} $\to$ \texttt{7.2.0} \\
\texttt{wujia\_portal\_debt} & \texttt{css/portal\_debt.css} (gỡ override) &
  \texttt{3.3.0} $\to$ \texttt{3.4.0} \\
\texttt{wujia\_portal\_support} & \texttt{views/portal\_support.xml} &
  \texttt{3.11.0} $\to$ \texttt{3.12.0} \\
\texttt{wujia\_portal\_return} & \texttt{views/portal\_return\_form.xml} &
  \texttt{2.3.0} $\to$ \texttt{2.4.0} \\
\hline
\end{tabular}
\end{center}

\textbf{Không} model/field mới, \textbf{không} migration, \textbf{không} đổi quyền hay dữ liệu.
Commit \texttt{70f3e4f} trên nhánh \texttt{dev/2026-08-11-b4}. \textbf{Chưa deploy.}

\section{Bài học}

\begin{enumerate}
\item \textbf{Blast radius lớn thì tìm seam, đừng sửa 31 chỗ.} Một helper trên
  \texttt{ir.http} gọi từ QWeb cho toàn bộ 31 call site hưởng hành vi mới với 0 dòng
  controller --- và cũng có nghĩa 0 chỗ để quên.
\item \textbf{Validate redirect phải nghĩ tới bước chuẩn hoá của trình duyệt.} Kiểm tiền tố
  \texttt{/portal/} là chưa đủ: \texttt{/portal/../web} qua được. Luật đúng là cấm segment
  \texttt{..} \emph{trước} khi xét tiền tố.
\item \textbf{Hai spec BA lệch nhau thì hỏi, đừng chọn.} 800 vs 700 cho title mobile: chọn 800
  sẽ phá bảng đo đã Pass của B3a. Ghi vào LIMIT và chuyển BA.
\item \textbf{Lỗi có sẵn thì miễn trừ tường minh, đừng nới ngưỡng} (L7). Ghi rõ trang nào,
  thông điệp nào, và bằng chứng nó không do sprint này (file 0 dòng diff).
\item \textbf{Harness sai thì sửa harness} (L7/L9). Trong sprint này riêng harness sai 6 lần
  (login sai route, user không đọc được ticket, đọc computed style thay vì pseudo-element,
  query lồng chưa encode, đọc URL trước khi điều hướng xong, selector không scope, chọn user
  mà role vốn không đổi, sai selector badge giỏ hàng) --- chân lý luôn là spec BA.
\item \textbf{\texttt{pkill} chung lệnh với lệnh khởi động thì shell tự giết mình} (L8) ---
  tách hai lần gọi.
\end{enumerate}

\section{LIMIT sau B4}

\begin{enumerate}
\item \textbf{Câu hỏi gửi BA} --- \texttt{CMP-BPH-001} ghi title mobile 800,
  \texttt{CMP-PG-001} ghi 700. Dev giữ 700, chờ BA chốt.
\item Cảnh báo chưa lưu hiện áp cho hai màn tạo mới; màn khác chưa có ô nhập nên chưa cần.
\item 21/23 component còn lại chưa có spec, gồm \texttt{CMP-DS-001} DetailSummary --- mã
  chứng từ vẫn nằm trong card thông tin sẵn có.
\item Hệ \texttt{.wj-pc-page-header} vẫn còn ở \texttt{/portal/exam/register},
  \texttt{/portal/exam/\{id\}}, \texttt{pc\_preview}.
\item Padding ngang PC vẫn do content-wrapper của shell Vuexy cấp --- chờ CMP PageContainer.
\item \textbf{Việc mới phát hiện}: \texttt{portal\_info\_request\_form.xml:122} có inline
  script chết vì \texttt{\&\&} bị escape --- lỗi có sẵn, ngoài scope B4, cần mở issue riêng.
\end{enumerate}
